\documentclass[../main.tex]{subfiles}
\begin{document}

% --- Local helpers (guarded so this subfile compiles even if the master
%     preamble does not already load these packages). Placed at the top of the
%     document body so the llmbox/\hi macros are defined before their first use
%     in the pipeline-ablation section. ---------------------------------------
\usepackage{tcolorbox}
\tcbuselibrary{breakable,skins}
\usepackage{soul}
\providecommand{\hi}[1]{\hl{#1}}            % inline highlight of the analysed span
\providecommand{\enquote}[1]{``#1''}        % fallback if csquotes is not loaded
% LLM reasoning-chain box
\newtcolorbox{llmbox}[1][]{
  breakable, enhanced, sharp corners,
  colback=black!3, colframe=black!55, boxrule=0.4pt,
  fonttitle=\bfseries\small, fontupper=\small\ttfamily,
  left=4pt, right=4pt, top=3pt, bottom=3pt,
  title={LLM reasoning chain (verbatim)}, #1}

% =============================================================================
\section{Experimental Setup}
\label{sec:setup}
% =============================================================================

\subsection{Hardware and Software Environment}

All offline indexing steps (embedding, FAISS index construction, profiling) run
on a local workstation with an 11th Gen Intel Core i5-1135G7 CPU at 2.40\,GHz
(reported boost 2.42\,GHz), 8\,GB RAM, and an NVIDIA GeForce MX350 2\,GB GPU.
FAISS approximate nearest-neighbour search operates entirely in-memory on CPU;
no GPU is required for retrieval. GPT-4o-mini classification calls are issued
through the OpenAI Chat Completions API, so the inference hardware behind the
API endpoint is out of scope. The reproduced LLaMA-3-8B runs are executed on
Kaggle using two NVIDIA Tesla T4 GPUs (T4$\times$2).

\subsection{Language Model Endpoints}

All runs use the \texttt{gpt-4o-mini} endpoint as accessed in May 2026.

\begin{table}[htbp]
  \centering
  \caption{Language model endpoints used in the experiments.}
  \label{tab:endpoints}
  \small
  \begin{tabular}{lll}
    \toprule
    \textbf{Role} & \textbf{Model} & \textbf{Used for} \\
    \midrule
    Offline profiler (train, labelled)   & GPT-4o-mini & Build training profiles with labels \\
    Offline profiler (test, label-blind) & GPT-4o-mini & Build test profiles without labels \\
    Classifier                           & GPT-4o-mini & All five prompt modes \\
    \bottomrule
  \end{tabular}
\end{table}

Decoding is deterministic in all runs: temperature is fixed at 0.0 and the
random seed is held constant across experiments.
The maximum output budget differs by prompt mode: compact modes (zero-shot,
one-shot) use a small token budget, while the Reasoned-RAG mode requires a
larger budget to accommodate the full XML reasoning chain.

\subsection{Reproducibility}

Every prediction run writes a timestamped log directory containing the full
prompt--response trail.
Reasoned-RAG runs additionally emit a structured reasoning log recording each
intermediate XML block produced during inference.
The run identifier is derived from the wall-clock timestamp at run start, so
every experiment is traceable to a specific log directory.

% =============================================================================
\section{Datasets}
\label{sec:datasets}
% =============================================================================

\subsection{Essays Corpus}

The primary evaluation corpus is the stream-of-consciousness essay dataset
collected by Pennebaker and King~\cite{pennebaker1999linguistic}, comprising
2{,}468 short essays written by undergraduate students.
Each essay is accompanied by five binary (high/low) Big Five personality
labels derived from a median split of self-report questionnaire scores.
The median-split procedure makes each trait approximately balanced in the
full dataset, with mild class drift in held-out splits.

The corpus is partitioned into three non-overlapping splits:
a training set of 1{,}974 essays used to build the retrieval index and supply
exemplars; a validation set of 247 essays used for retrieval-strategy
ablation and gamma-weight optimisation; and a test set of 247 essays used
for final classification evaluation.

\subsection{Dataset Statistics}

\begin{table}[htbp]
  \centering
  \caption{Essays dataset statistics. Labels are derived from a median split
    of self-report questionnaire scores.
    \%High = fraction of samples labelled \textit{high} for that trait.}
  \label{tab:datasets}
  \small
  \begin{tabular}{lrccccc}
    \toprule
    \textbf{Split} & \textbf{$n$}
      & \textbf{OPN} & \textbf{CON} & \textbf{EXT}
      & \textbf{AGR} & \textbf{NEU} \\
    \midrule
    Train & 1{,}974 & 0.515 & 0.510 & 0.515 & 0.515 & 0.510 \\
    Val   &   247   & 0.530 & 0.482 & 0.524 & 0.486 & 0.514 \\
    Test  &   247   & 0.508 & 0.530 & 0.518 & 0.567 & 0.498 \\
    \bottomrule
  \end{tabular}
\end{table}

% =============================================================================
\section{Evaluation Metrics}
\label{sec:metrics}
% =============================================================================

\subsection{Classification Metrics}

All classification metrics are computed independently per trait, then summarised
across the five OCEAN traits.
Let $\mathrm{TP}_c$, $\mathrm{FP}_c$, $\mathrm{FN}_c$ denote true positives,
false positives, and false negatives for class $c \in \{\textit{high},
\textit{low}\}$.
Per-class precision and recall are:
%
\begin{equation}
  P_c = \frac{\mathrm{TP}_c}{\mathrm{TP}_c + \mathrm{FP}_c},\qquad
  R_c = \frac{\mathrm{TP}_c}{\mathrm{TP}_c + \mathrm{FN}_c}.
\end{equation}
%
The per-class F1 score is the harmonic mean of precision and recall:
%
\begin{equation}
  F1_c = \frac{2 P_c R_c}{P_c + R_c}.
\end{equation}
%
Macro-F1 averages F1 equally across the two classes, giving equal weight to
the minority class:
%
\begin{equation}
  \text{Macro-F1} = \frac{1}{2}\left(F1_{\textit{high}} + F1_{\textit{low}}\right).
\end{equation}
%
Macro-F1 is the headline metric throughout this chapter because the test splits
exhibit mild class imbalance per trait.

\subsection{Retrieval Metrics}

Retrieval quality is assessed on the validation split ($n=247$) to avoid
contaminating the test labels used for classification evaluation.
Two metrics are reported per (strategy, trait, $k$) combination:

\begin{description}
  \item[Mean-match-rate@$k$ (MMR)] The average fraction of the top-$k$
    retrieved training essays whose ground-truth label for the target trait
    matches that of the query essay:
    \begin{equation}
      \text{MMR@}k = \frac{1}{|Q|} \sum_{q \in Q}
        \frac{\bigl|\{i \leq k : \ell_i = \ell_q\}\bigr|}{k},
    \end{equation}
    where $Q$ is the set of queries and $\ell_q$ is the query's true label.
    A value near 0.5 corresponds to chance; higher is better.

  \item[Hit-rate@$k$ (HR)] The fraction of queries for which at least one of
    the top-$k$ retrieved essays has a matching label:
    \begin{equation}
      \text{HR@}k = \frac{1}{|Q|}
        \sum_{q \in Q} \mathbf{1}\!\left[\exists\, i \leq k : \ell_i = \ell_q\right].
    \end{equation}
    Hit-rate rewards diversity in the retrieved pool rather than overall
    label purity.
\end{description}

Both metrics are computed per trait and then macro-averaged across the five
OCEAN traits for aggregate comparison.

% =============================================================================
\section{Baseline Comparison}
\label{sec:baselines}
% =============================================================================

To place the proposed pipeline in context, this section reports a set of
baselines on the Essays test set. The classical and neural baselines (SVM,
LSTM, BERT, RoBERTa, DeBERTa) are supervised classifiers trained on the Essays
training split; the LLM baselines (GPT-4o-mini and LLaMA-3-8B under zero-shot,
one-shot, and chain-of-thought prompting) are non-retrieval prompting
baselines.

A note on the purpose of this comparison: the fine-tuned encoder baselines
(BERT, RoBERTa, DeBERTa) are directly optimised for classification accuracy on
this benchmark and have access to full supervision over all training examples.
The proposed pipeline is designed with a different objective --- to produce an
interpretable, structured psychological representation that disentangles surface
writing style from trait-relevant evidence, and to expose the reasoning process
to inspection via 30-facet profiles and structured XML chains.
These design choices impose a cost in raw discriminative power that the
comparison below makes visible. The baseline table is therefore best read as
a diagnostic of that cost rather than as a competitive ranking.
All figures are macro-averaged F1 per trait (the mean of the high-class and
low-class F1), and the rightmost column is the average over the five traits.
The PADO row is quoted from the original publication~\cite{yeo2025pado} and uses
GPT-3.5; it is included only as an external reference point and was not
reproduced here.
The proposed Reasoned-RAG pipeline (described fully in
Section~\ref{sec:pipeline-ablation}) is included for direct comparison;
those rows correspond to the same runs reported in
Tables~\ref{tab:pipeline-ablation} and~\ref{tab:model-comparison}.

\subsection{Baseline Configurations}
\label{sec:baseline-config}

All reproduced baselines are trained and evaluated per trait (one independent
binary classifier per OCEAN dimension) on the Essays splits
(train $n=1{,}974$, validation $n=247$, test $n=247$), with the validation
split used for model selection. Their settings are as follows.

\paragraph{SVM.}
A support vector machine with an RBF kernel ($C=10$, $\gamma=0.1$) over TF--IDF
features, trained independently per trait. The TF--IDF vectoriser uses word uni-
and bi-grams ($\text{ngram\_range}=(1,2)$), a minimum document frequency of $2$,
sublinear term-frequency scaling, and a vocabulary capped at $10{,}000$
features. A fixed random seed ($42$) is used.

\paragraph{LSTM.}
A bidirectional LSTM with a shared text encoder and one classification head per
trait. Text is lower-cased and reduced to alphanumeric word tokens; a
vocabulary is built from the training split with a minimum token frequency of
$2$, and sequences are truncated or padded to $512$ tokens. The model uses a
$128$-dimensional embedding, two stacked bidirectional LSTM layers with hidden
size $128$ per direction, dropout $0.3$, mask-aware mean-pooling over the hidden
states, and a $64$-unit projection before the per-trait sigmoid heads.
Training uses Adam (learning rate $10^{-3}$), batch size $32$, binary
cross-entropy loss, gradient clipping at $1.0$, up to $50$ epochs with early
stopping (patience $4$ on average validation F1), and a decision threshold of
$0.5$.

\paragraph{Transformer encoders (BERT, RoBERTa, DeBERTa).}
Pre-trained base-size encoders (BERT-base-uncased, RoBERTa-base,
DeBERTa-v3-base) fine-tuned with a shared encoder and one linear head per
trait on the \texttt{[CLS]} representation, using identical hyper-parameters
across the three backbones. Inputs are tokenised with each model's tokeniser and truncated
or padded to $512$ tokens. Fine-tuning uses AdamW with layer-wise
learning-rate decay (base learning rate $2\times10^{-5}$, decay factor $0.95$
per layer, head learning rate $10\times$ the base), weight decay $0.01$, a
linear schedule with $10\%$ warm-up, batch size $16$, binary cross-entropy
loss, gradient clipping at $1.0$, and up to $10$ epochs
with early stopping (patience $3$ on average validation F1).

\paragraph{LLM prompting baselines.}
GPT-4o-mini and LLaMA-3-8B-Instruct are evaluated under zero-shot, one-shot,
and chain-of-thought prompting without retrieval, using the prompt templates
described in the Methodology chapter and deterministic decoding (temperature
$0$).

\begin{table}[htbp]
  \centering
  \caption{Macro-F1 per trait on the Essays test set. The PADO row is
    reported from the literature (GPT-3.5) and is not directly comparable to the
    locally reproduced rows. Best reproduced average in \textbf{bold}.}
  \label{tab:baselines}
  \small
  \begin{tabular}{llcccccr}
    \toprule
    \textbf{Model} & \textbf{Setting} & \textbf{OPN} & \textbf{CON}
      & \textbf{EXT} & \textbf{AGR} & \textbf{NEU} & \textbf{Avg} \\
    \midrule
    PADO (GPT-3.5)\textsuperscript{$\dagger$} & ---       & 0.720 & 0.680 & 0.690 & 0.670 & 0.670 & 0.686 \\
    \midrule
    SVM        & ---       & 0.563 & 0.538 & 0.521 & 0.528 & 0.577 & 0.545 \\
    LSTM       & ---       & 0.558 & 0.515 & 0.542 & 0.537 & 0.511 & 0.533 \\
    BERT       & ---       & 0.609 & 0.540 & 0.533 & 0.568 & 0.599 & 0.570 \\
    RoBERTa    & ---       & 0.624 & 0.534 & 0.617 & 0.515 & 0.524 & 0.563 \\
    DeBERTa    & ---       & 0.570 & 0.548 & 0.546 & 0.484 & 0.496 & 0.529 \\
    \midrule
    GPT-4o-mini & zero-shot & 0.585 & 0.440 & 0.495 & 0.555 & 0.450 & 0.505 \\
                & one-shot  & 0.595 & 0.525 & 0.515 & 0.560 & 0.550 & 0.549 \\
                & CoT       & 0.480 & 0.460 & 0.510 & 0.555 & 0.450 & 0.491 \\
    \midrule
    LLaMA-3-8B  & zero-shot & 0.520 & 0.470 & 0.540 & 0.505 & 0.495 & 0.506 \\
                & one-shot  & 0.530 & 0.530 & 0.510 & 0.565 & 0.555 & 0.538 \\
                & CoT       & 0.505 & 0.450 & 0.555 & 0.520 & 0.480 & 0.502 \\
    \midrule
    \textbf{Ours (GPT-4o-mini)}\textsuperscript{$\ddagger$} & Reasoned-RAG & \textbf{0.587} & \textbf{0.486} & \textbf{0.522} & \textbf{0.550} & 0.409 & \textbf{0.511} \\
    \textbf{Ours (LLaMA-3-8B)}\textsuperscript{$\ddagger$}  & Reasoned-RAG & 0.452 & 0.499 & 0.481 & 0.431 & \textbf{0.527} & 0.478 \\
    \bottomrule
  \end{tabular}
\end{table}

\noindent\textsuperscript{$\dagger$}\,Reported from~\cite{yeo2025pado}; not reproduced
in this work.
\noindent\textsuperscript{$\ddagger$}\,These rows are the same runs reported in
Tables~\ref{tab:pipeline-ablation} and~\ref{tab:model-comparison}.

Among the locally reproduced baselines, the fine-tuned transformer encoders are
the strongest, with BERT achieving the highest average macro-F1 ($0.570$). The
classical SVM ($0.545$) is competitive with the larger neural models, and the
non-retrieval LLM prompting baselines fall in a similar range ($0.49$--$0.55$),
with one-shot prompting outperforming both zero-shot and chain-of-thought for
both GPT-4o-mini and LLaMA-3-8B. The proposed Reasoned-RAG pipeline with
GPT-4o-mini ($0.511$) falls below BERT but is competitive with the non-retrieval
LLM baselines; with LLaMA-3-8B ($0.478$) the pipeline performs below the GPT-4o-mini
variant on most traits. The externally reported PADO figure ($0.686$) is
substantially higher, but it uses a different (larger) model and was not
reproduced under the present protocol, so it should not be read as a controlled
comparison.

% =============================================================================
\section{Retrieval Ablation}
\label{sec:retrieval-ablation}
% =============================================================================

% HOW TO RUN:
%   notebook/ablation/rag_ablation_rawpost.ipynb      -> rawpost
%   notebook/ablation/rag_ablation_profile.ipynb      -> profile
%   notebook/ablation/[FIX] rag_ablation_sliced_dual.ipynb -> sliced_dual
%   notebook/ablation/rag_ablation_facet_embed.ipynb  -> facet_embed
%   notebook/ablation/rag_ablation_hybrid_facet.ipynb -> hybrid_facet (gamma sweep)
%   Summaries saved under result/retrieve_accuracy/<strategy>/accuracy_summary.csv
%   and result/retrieve_accuracy/hybrid_facet/best_gamma.csv.

This section compares five retrieval strategies across two values of $k$
($k \in \{3,\,5\}$) on the Essays validation set.
Each strategy differs in how the query embedding is constructed and which
training index is searched.
The five strategies are described in Table~\ref{tab:strategies}.

\begin{table}[htbp]
  \centering
  \caption{Retrieval strategies compared in the ablation.}
  \label{tab:strategies}
  \small
  \begin{tabular}{lp{9cm}}
    \toprule
    \textbf{Strategy} & \textbf{Description} \\
    \midrule
    \textit{Raw-post}     & Embeds the raw essay text; searches a raw-text FAISS index.
                            Serves as the surface-style baseline. \\
    \textit{Profile}      & Embeds the full 30-facet profile text serialisation;
                            searches a profile-embedding FAISS index.
                            Removes lexical confound but ignores raw text. \\
    \textit{Sliced-dual}  & Fuses a raw-text embedding with a \emph{trait-sliced}
                            6-facet profile embedding; searches the full dual index.
                            The query embedding is trait-specific. \\
    \textit{Facet-embed}  & Fuses a raw-text embedding with the full 30-facet profile
                            embedding (dual fusion, $\gamma=0$ baseline of hybrid). \\
    \textit{Hybrid-facet} & Extends Facet-embed with a per-trait trait-masked facet
                            cosine re-ranking step, combining dense and structured
                            signals via Equation~\eqref{eq:hybrid}.
                            This is the proposed method. \\
    \bottomrule
  \end{tabular}
\end{table}

\subsection{Gamma Selection for Hybrid-Facet Retrieval}

The hybrid score in Equation~\eqref{eq:hybrid} is controlled by a per-trait
weight $\gamma_\tau$ that governs the contribution of the structured facet
channel relative to the dense channel.
A sweep over $\gamma \in \{0.0,\,0.2,\,0.4,\,0.6,\,0.8,\,1.0\}$ is conducted
on the validation set, with $\beta = 1-\gamma$ throughout.
The best $\gamma$ is selected per trait by maximising mean-match-rate@$k$.
Table~\ref{tab:best-gamma} reports the selected weights.

\begin{table}[htbp]
  \centering
  \caption{Best $\gamma$ per trait selected by validation MMR. These weights and
    MMR values come from a dedicated gamma-selection sweep over the validation
    set; they are reported separately from the cross-strategy comparison in
    Table~\ref{tab:retrieval-ablation}, which uses a single later run with a
    fixed per-trait $\gamma$, so the hybrid-facet MMR values in the two tables
    are not expected to coincide exactly.}
  \label{tab:best-gamma}
  \small
  \begin{tabular}{lcccc}
    \toprule
    \textbf{Trait} & \multicolumn{2}{c}{$k=3$} & \multicolumn{2}{c}{$k=5$} \\
    \cmidrule(lr){2-3}\cmidrule(lr){4-5}
     & $\gamma^*$ & MMR & $\gamma^*$ & MMR \\
    \midrule
    Openness          & 1.0 & 0.580 & 1.0 & 0.568 \\
    Conscientiousness & 0.8 & 0.572 & 1.0 & 0.573 \\
    Extraversion      & 0.2 & 0.571 & 0.2 & 0.565 \\
    Agreeableness     & 0.2 & 0.572 & 0.4 & 0.566 \\
    Neuroticism       & 1.0 & 0.596 & 1.0 & 0.585 \\
    \bottomrule
  \end{tabular}
\end{table}

\subsection{Retrieval Results}

Table~\ref{tab:retrieval-ablation} compares MMR and HR across all five
strategies at $k \in \{3,\,5\}$.
The hybrid-facet strategy uses the per-trait $\gamma^*$ values from
Table~\ref{tab:best-gamma}.

\begin{table}[htbp]
  \centering
  \caption{Retrieval ablation: mean-match-rate (MMR) and hit-rate (HR)
    at $k \in \{3,\,5\}$ across five strategies on the Essays validation set
    ($n=247$). Each cell shows MMR\,/\,HR. Best MMR per column in \textbf{bold}.}
  \label{tab:retrieval-ablation}
  \small
  \begin{tabular}{llccccc}
    \toprule
    \textbf{Strategy} & $k$
      & \textbf{OPN} & \textbf{CON} & \textbf{EXT}
      & \textbf{AGR} & \textbf{NEU} \\
    \midrule
    \multirow{2}{*}{\textit{Raw-post}}
      & 3 & 0.494\,/\,0.870 & 0.501\,/\,0.891 & 0.491\,/\,0.870 & 0.507\,/\,0.838 & 0.548\,/\,0.887 \\
      & 5 & 0.504\,/\,0.951 & 0.504\,/\,0.984 & 0.504\,/\,0.972 & 0.512\,/\,0.960 & 0.538\,/\,0.968 \\
    \addlinespace
    \multirow{2}{*}{\textit{Profile}}
      & 3 & 0.514\,/\,0.878 & 0.526\,/\,0.903 & 0.526\,/\,0.878 & 0.518\,/\,0.826 & \textbf{0.560}\,/\,0.846 \\
      & 5 & 0.520\,/\,0.960 & 0.511\,/\,0.960 & 0.509\,/\,0.943 & \textbf{0.534}\,/\,0.927 & \textbf{0.554}\,/\,0.939 \\
    \addlinespace
    \multirow{2}{*}{\textit{Sliced-dual}}
      & 3 & 0.526\,/\,0.854 & 0.529\,/\,0.854 & \textbf{0.549}\,/\,0.875 & \textbf{0.525}\,/\,0.838 & 0.547\,/\,0.842 \\
      & 5 & \textbf{0.538}\,/\,0.951 & \textbf{0.532}\,/\,0.951 & \textbf{0.546}\,/\,0.947 & 0.512\,/\,0.939 & 0.552\,/\,0.939 \\
    \addlinespace
    \multirow{2}{*}{\textit{Facet-embed}}
      & 3 & \textbf{0.533}\,/\,0.878 & \textbf{0.538}\,/\,0.903 & 0.536\,/\,0.895 & 0.487\,/\,0.830 & 0.528\,/\,0.883 \\
      & 5 & 0.536\,/\,0.968 & 0.524\,/\,0.976 & 0.537\,/\,0.988 & 0.493\,/\,0.943 & 0.540\,/\,0.980 \\
    \addlinespace
    \multirow{2}{*}{\textbf{\textit{Hybrid-facet}}\textsuperscript{$\S$}}
      & 3 & \textbf{0.533}\,/\,0.878 & \textbf{0.538}\,/\,0.903 & 0.536\,/\,0.895 & 0.487\,/\,0.830 & 0.528\,/\,0.883 \\
      & 5 & 0.536\,/\,0.968 & 0.524\,/\,0.976 & 0.537\,/\,0.988 & 0.493\,/\,0.943 & 0.540\,/\,0.980 \\
    \bottomrule
  \end{tabular}

\noindent\textsuperscript{$\S$}\,The per-trait facet re-ranking did not alter the
top-$k$ relative to Facet-embed in this run; rows are therefore identical.
\end{table}

Across all five strategies the mean-match-rate stays close to the chance level
of 0.5, which reflects the difficulty of label-pure retrieval on the Essays
corpus: neighbours that are similar in style or topic frequently carry the
opposite trait label. The differences between strategies are therefore small
(within roughly $0.05$ MMR) and no single strategy dominates across all traits.
The profile-embedding and the trait-sliced dual strategies give the best MMR on
most traits---profile is best on Neuroticism at both $k$ values ($0.560$ at
$k=3$) and on Agreeableness at $k=5$, while sliced-dual is best on Extraversion
and Agreeableness at $k=3$ and on Openness, Conscientiousness and Extraversion
at $k=5$. The facet-embed strategy leads on Openness and Conscientiousness at
$k=3$. The raw-post baseline is the weakest on most traits, consistent with the
expectation that surface style is a poor proxy for personality labels.

In this run the per-trait facet re-ranking of the hybrid-facet strategy did not
change the selected top-$k$ relative to facet-embed, so the two rows are
identical in Table~\ref{tab:retrieval-ablation}; the structured re-ranking
therefore provides no measurable retrieval gain over plain facet fusion here.
Hit-rate is high for every strategy ($0.83$--$0.99$) and increases with $k$ as
expected, so the small MMR differences are not driven by candidate-pool
coverage. Overall, the retrieval ablation shows that incorporating the
structured facet signal (profile, sliced-dual, facet-embed) yields a modest but
consistent MMR improvement over the raw-text baseline, while the additional
trait-masked re-ranking step does not improve retrieval on this validation set.

% =============================================================================
\section{Pipeline Component Ablation}
\label{sec:pipeline-ablation}
% =============================================================================

% HOW TO RUN:
%   Full system : notebook/gpt/rag_profile_half2_predict copy.ipynb (reasoned_rag_def_oneshot_30f)
%   No-RAG      : notebook/ablation_main_pipe/ablation_no_rag.ipynb
%   No-reasoning: notebook/ablation_main_pipe/ablation_no_reasoning.ipynb
%   Each saves predictions.csv + evaluation_summary.csv; macro_f1 column used below.

This section ablates the two principal components of the proposed pipeline,
retrieval and structured reasoning, while holding the rest of the configuration
fixed (GPT-4o-mini, Essays test set $n=247$). The full system corresponds to
Reasoned-RAG-Def-Oneshot. The \emph{no-RAG} variant removes the retrieved
exemplars and keeps the structured reasoning prompt; the \emph{no-reasoning}
variant keeps the retrieved 30-facet exemplars but replaces the XML reasoning
chain with a direct label.

\begin{table}[htbp]
  \centering
  \caption{Component ablation on the Essays test set ($n=247$, GPT-4o-mini).
    Per-trait and averaged macro-F1. Best per column in \textbf{bold}.}
  \label{tab:pipeline-ablation}
  \small
  \begin{tabular}{lcccccr}
    \toprule
    \textbf{Configuration} & \textbf{OPN} & \textbf{CON} & \textbf{EXT}
      & \textbf{AGR} & \textbf{NEU} & \textbf{Macro} \\
    \midrule
    \textbf{Full system} (Reasoned-RAG)  & 0.587 & 0.486 & 0.522 & 0.550 & 0.409 & 0.511 \\
    \quad $-$ retrieval (no-RAG)          & 0.540 & 0.452 & 0.464 & 0.514 & 0.357 & 0.465 \\
    \quad $-$ reasoning (no-reasoning)    & \textbf{0.597} & 0.398 & \textbf{0.554} & \textbf{0.537} & \textbf{0.553} & \textbf{0.528} \\
    \bottomrule
  \end{tabular}
\end{table}

Removing retrieval lowers the average macro-F1 from $0.511$ to $0.465$, and the
drop is present on every trait, indicating that the retrieved exemplars
contribute positively to the full system across the board.
Removing the structured reasoning chain while retaining the retrieved exemplars
raises the average macro-F1 to $0.528$, ahead of the full system on four
of five traits (all except Conscientiousness). Taken at face value this would
suggest the XML reasoning chain is unhelpful; the per-class analysis in
Section~\ref{sec:ablation-perclass} shows the picture is more nuanced, and traces
the macro difference to the anti-bias floor and its interaction with the
state/trait distinction rather than to a uniform effect of the chain.

\subsection{Per-Class Breakdown of the Ablation}
\label{sec:ablation-perclass}

The macro-F1 figures in Table~\ref{tab:pipeline-ablation} hide \emph{why} the
no-reasoning variant wins. Because macro-F1 averages the \textit{high}- and
\textit{low}-class F1, an equal macro can arise either from balanced per-class
performance or from a strong skew toward one pole. Table~\ref{tab:ablation-perclass}
therefore decomposes each configuration into its per-class (high\,/\,low) F1
and reports the directed error counts ($\text{gold=high}\rightarrow\text{pred=low}$
versus the reverse) that drive the skew.

\begin{table}[htbp]
  \centering
  \caption{Per-class breakdown of the component ablation
    (Essays test set, $n=247$, GPT-4o-mini). For each trait the table reports
    \textit{high}-class F1, \textit{low}-class F1, and the two directed error
    counts $H{\to}L$ (gold \textit{high} predicted \textit{low}) and
    $L{\to}H$ (the reverse). A large $H{\to}L$ relative to $L{\to}H$ indicates
    suppression of the \textit{high} pole; the reverse indicates over-prediction
    of \textit{high}.}
  \label{tab:ablation-perclass}
  \small
  \begin{tabular}{llccccc}
    \toprule
    \textbf{Config} & \textbf{Metric}
      & \textbf{OPN} & \textbf{CON} & \textbf{EXT}
      & \textbf{AGR} & \textbf{NEU} \\
    \midrule
    \multirow{4}{*}{\textbf{Full system}}
      & F1 (high)     & 0.585 & 0.328 & 0.358 & 0.528 & 0.166 \\
      & F1 (low)      & 0.589 & 0.644 & 0.687 & 0.571 & 0.653 \\
      & $H{\to}L$     & 55    & 97    & 98    & 71    & 112   \\
      & $L{\to}H$     & 47    & 18    & 6     & 40    & 9     \\
    \addlinespace
    \multirow{4}{*}{$-$ retrieval}
      & F1 (high)     & 0.438 & 0.238 & 0.253 & 0.426 & 0.047 \\
      & F1 (low)      & 0.642 & 0.665 & 0.674 & 0.603 & 0.667 \\
      & $H{\to}L$     & 85    & 107   & 108   & 90    & 121   \\
      & $L{\to}H$     & 23    & 8     & 4     & 26    & 1     \\
    \addlinespace
    \multirow{4}{*}{$-$ reasoning}
      & F1 (high)     & 0.568 & 0.120 & 0.491 & 0.459 & 0.683 \\
      & F1 (low)      & 0.626 & 0.676 & 0.617 & 0.616 & 0.423 \\
      & $H{\to}L$     & 62    & 117   & 75    & 86    & 15    \\
      & $L{\to}H$     & 37    & 0     & 33    & 25    & 86    \\
    \bottomrule
  \end{tabular}
\end{table}

The headline reading of Table~\ref{tab:pipeline-ablation} is that removing
reasoning is, on average, slightly better than the full system ($0.528$ vs.\
$0.511$), but the per-class view in Table~\ref{tab:ablation-perclass} shows this
average is misleading: on three of five traits (OPN, EXT, AGR) the two
configurations are in fact \emph{close}, and the entire macro gap is produced by
\emph{two traits that move in opposite directions}. On Conscientiousness the full
system is much stronger (macro $0.486$ vs.\ $0.398$), while on Neuroticism the
no-reasoning variant is much stronger (macro $0.409$ vs.\ $0.553$). These two
traits, rather than the three near-ties, are where the reasoning chain earns or
loses its keep, so the analysis below focuses on them.

Crucially, both gaps are driven by the \emph{same} state/trait gate operating in
opposite regimes. On Conscientiousness the gate works \emph{for} the chain: the
no-reasoning variant collapses almost entirely to \textit{low} (high-class F1
$0.120$, with $117$ $H{\to}L$ errors and \emph{zero} $L{\to}H$), whereas the
chain's explicit search for trait-level discipline cues recovers a sizeable share
of the true \textit{high} cases (high-class F1 $0.328$). On Neuroticism the gate
works \emph{against} the chain: anxiety and worry in these essays are almost
always phrased as momentary states, so the gate tags them \texttt{[state]} and
forces \textit{low} ($112$ $H{\to}L$ errors, high-class F1 $0.166$); removing the
gate lets the bare-label model read the surface stress vocabulary directly and
flips the bias the other way ($86$ $L{\to}H$ errors, high-class F1 up to $0.683$
but low-class F1 down to $0.423$). The same mechanism that rescues
Conscientiousness therefore sinks Neuroticism. The two worked examples below show
each regime concretely.

\paragraph{Conscientiousness: the gate rescues a true \textit{high}.}
Record~48 (Conscientiousness, gold \textit{high}) reports concrete, habitual
self-discipline---``I make my bed everyday\dots I get all my assignments done
early\dots I finished my scholarship applications.'' These are exactly the
recurring, cross-situational behaviours the gate is designed to credit, so the
chain tags them \texttt{[trait]} and commits to \textit{high}; the bare-label
variant, with no such structured search, returns \textit{low}:
%
\begin{llmbox}[title={Record 48 -- Conscientiousness (gold: high)}]
\textbf{Full system (reasoned-RAG) $\rightarrow$ predicts \hi{high}} ~~(correct)\\
\textit{Evidence.} \hi{[trait]} ``I make my bed everyday, I get all my assignments done early\dots I finished my scholarship applications.''\\
\textit{Facet check.} competence/order/dutifulness/achievement/self-discipline \hi{all $\to$ high}.\\
\textit{Cue tally.} [trait] cues: \hi{4} | [state] cues: 0 ~~$\Rightarrow$~ Presumption: HIGH.\\
\textit{Verdict.} ``consistent evidence of organization, responsibility, and self-discipline, leading to a conclusion of high.''\\[4pt]
\textbf{$-$ reasoning $\rightarrow$ predicts \hi{low}} ~~(wrong)
\end{llmbox}
%
This is why Conscientiousness is the one trait that \emph{needs} the chain:
when the high-pole cues are genuinely dispositional, the state/trait gate
correctly admits them, and without the gate the model defaults toward
\textit{low} (zero $L{\to}H$ errors in the no-reasoning column). Record~48 is one
of the high-class cases the chain recovers and the bare-label model loses.

\paragraph{Neuroticism: the gate discards a correct \textit{low} once removed.}
Neuroticism is the mirror image. Record~6 (gold \textit{low}) describes ordinary
pre-exam nerves---``I feel a bit stressed\dots my first three tests are all next
week.'' Here the gate is \emph{right} to discount the cue: the worry is
situational, the gold label is \textit{low}, and the chain reaches \textit{low}.
Removing the chain destroys exactly this distinction:
%
\begin{llmbox}[title={Record 6 -- Neuroticism (gold: low)}]
\textbf{Full system (reasoned-RAG) $\rightarrow$ predicts \hi{low}} ~~(correct)\\
\textit{Facet check.} anxiety \hi{$\to$ high} (dominant evidence type: \hi{state}); vulnerability \hi{$\to$ high (state)}; others low.\\
\textit{Cue tally.} [trait] cues: \hi{0} | [state] cues: 3 ~~$\Rightarrow$~ Presumption: LOW.\\
\textit{Verdict.} ``situational stress and worry about tests but lacks any recurring, trait-level evidence\dots classified as low.''\\[4pt]
\textbf{$-$ reasoning $\rightarrow$ predicts \hi{high}} ~~(wrong)
\end{llmbox}
%
The bare-label variant, lacking any state/trait distinction, takes the surface
stress vocabulary at face value and predicts \textit{high}---one of $86$ such
$L{\to}H$ Neuroticism errors that the removed chain had previously prevented.
Note that on this trait the chain's facet stage \emph{does} score anxiety and
vulnerability high; it is the \texttt{[state]} tag, not blindness to the cue,
that produces the correct \textit{low}---the same suppression mechanism examined
in full in Section~\ref{sec:qual-gpt}, here operating in its favour.

Taken together, the two records explain the whole CON/NEU divergence: the
state/trait gate is a single rule that helps when high-pole cues are genuinely
dispositional (Conscientiousness) and hurts when they are genuinely transient
yet the gold label still says \textit{high} (Neuroticism). Neither configuration
is well calibrated---the chain over-suppresses \textit{high}, deleting it
over-predicts \textit{high}---so the $+0.017$ macro advantage of the no-reasoning
variant is not evidence that reasoning is harmful in general, but that on this
corpus the Neuroticism loss happens to outweigh the Conscientiousness gain.
Calibrating the gate, rather than removing it, is therefore the appropriate fix.

% =============================================================================
\section{Classifier Backbone Comparison}
\label{sec:model-comparison}
% =============================================================================

The proposed pipeline was also run with an open-weight classifier,
LLaMA-3-8B-Instruct, in place of GPT-4o-mini, to assess whether the method
depends on a proprietary model. Both classifiers are evaluated on the full
test set ($n=247$). Table~\ref{tab:model-comparison} reports per-class F1
broken down by high and low poles together with the macro-F1 average; a
large gap between the two classes for a given trait indicates systematic
bias toward one pole.

\begin{table}[htbp]
  \centering
  \caption{Proposed pipeline (Reasoned-RAG) with two classifier backbones on
    the Essays test set ($n=247$). Per-class F1 (High\,/\,Low) and macro-F1
    per trait. A large High--Low gap indicates systematic bias toward one pole.}
  \label{tab:model-comparison}
  \small
  \begin{tabular}{llcccccr}
    \toprule
    \textbf{Classifier} & \textbf{Class}
      & \textbf{OPN} & \textbf{CON} & \textbf{EXT}
      & \textbf{AGR} & \textbf{NEU} & \textbf{Macro} \\
    \midrule
    \multirow{2}{*}{\mbox{GPT-4o-mini}}
      & High & 0.585 & 0.328 & 0.358 & 0.528 & 0.166 & \\
      & Low  & 0.589 & 0.644 & 0.687 & 0.571 & 0.653 & \\
    \cmidrule(lr){2-8}
      & \textit{Macro} & \textit{0.587} & \textit{0.486} & \textit{0.522} & \textit{0.550} & \textit{0.409} & \textit{0.511} \\
    \addlinespace
    \multirow{2}{*}{\mbox{LLaMA-3-8B}} % chktex 8
      & High & 0.629 & 0.523 & 0.522 & 0.564 & 0.606 & \\
      & Low  & 0.274 & 0.474 & 0.440 & 0.298 & 0.449 & \\
    \cmidrule(lr){2-8}
      & \textit{Macro} & \textit{0.452} & \textit{0.499} & \textit{0.481} & \textit{0.431} & \textit{0.527} & \textit{0.478} \\
    \bottomrule
  \end{tabular}
\end{table}

The two classifiers exhibit opposite bias directions.
GPT-4o-mini consistently under-predicts the \textit{high} pole: the gap
between low-class and high-class F1 is largest on Neuroticism ($0.653$ vs.\
$0.166$), Extraversion ($0.687$ vs.\ $0.358$), and Conscientiousness
($0.644$ vs.\ $0.328$), indicating that the anti-bias reasoning rule in the
XML chain over-corrects toward \textit{low} on these traits.
LLaMA-3-8B shows the reverse pattern, consistently over-predicting
\textit{high}: the gap is largest on Openness ($0.629$ vs.\ $0.274$) and
Agreeableness ($0.564$ vs.\ $0.298$).
Openness is the only trait where both classifiers achieve near-equal class
F1, suggesting it is the least affected by systematic bias under the current
prompting scheme.

% =============================================================================
\section{Qualitative Error Analysis}
\label{sec:qualitative}
% =============================================================================

This section presents a qualitative examination of representative prediction
errors, with a focus on the traits exhibiting the most severe systematic
bias identified in Table~\ref{tab:model-comparison}.

\subsection{Suppression of High Predictions in GPT-4o-mini}
\label{sec:qual-gpt}

All examples in this subsection are drawn from the same Reasoned-RAG run that
produced Table~\ref{tab:model-comparison}
(\texttt{reasoned\_rag\_def\_oneshot\_30f}, run \texttt{20260523-221029}).
The class-F1 gap in that table shows that GPT-4o-mini under-predicts the
\textit{high} pole on every trait; counting the directed errors in the run
confirms that the bias is overwhelmingly one-directional. The number of
\mbox{gold=high$\rightarrow$pred=low} errors dwarfs the reverse on the
worst-affected traits: Neuroticism ($112$ vs.\ $9$), Extraversion ($98$ vs.\
$6$), and Conscientiousness ($97$ vs.\ $18$), with Agreeableness ($71$ vs.\
$40$) and Openness ($55$ vs.\ $47$) less skewed but still net-suppressive.

Reading the reasoning chains for these false negatives reveals a single,
remarkably consistent mechanism rather than a scatter of unrelated mistakes.
The final \emph{verdict} stage of the XML chain applies an anti-bias rule that
demands \emph{stable, recurring, trait-level} evidence before it will commit to
\textit{high}, and discounts anything it has tagged as a momentary
\emph{state}. Because the Essays are short, in-the-moment diary entries, almost
no cue survives this filter, so the verdict collapses to \textit{low} even when
the earlier facet stage has already recorded high-pole evidence. The three
failure modes below isolate distinct symptoms of this same over-correction.

\paragraph{The verdict overrides high-scored facets.}
The clearest evidence that the bias lives in the verdict stage---not in the
evidence-gathering---is that the model frequently scores facets \textit{high}
and then concludes \textit{low}. In Neuroticism record~0 (gold \textit{high},
predicted \textit{low}) the facet check scores two of six facets at the high
pole, yet the verdict dismisses them:
%
\begin{llmbox}
\textbf{Facet check.}\\
- anxiety -> \hi{high} | expresses feeling ``stressed out'' and overwhelmed by school responsibilities (dominant evidence type: state)\\
- hostility -> low | no evidence of anger or resentment is present.\\
- depression -> none | no signs of sadness or hopelessness.\\
- self-conscious -> mod | reflects on feelings about being in a sorority.\\
- impulsiveness -> none | no clear evidence.\\
- vulnerability -> \hi{high} | expresses feeling overwhelmed and needing to manage priorities.\\[4pt]
\textbf{Verdict.}\\
The text indicates \hi{situational} stress and anxiety but \hi{lacks any stable, recurring patterns} of emotional instability characteristic of high Neuroticism. Therefore, it is classified as low.
\end{llmbox}
%
The two highlighted facets (\texttt{anxiety} and \texttt{vulnerability}) are
exactly the cues a human annotator would read as high Neuroticism, and the
ground-truth label agrees. The verdict does not dispute the cues; it
\emph{reclassifies them as situational} and applies the
\enquote{stable, recurring} gate, which the cues cannot pass because the essay
is a single time-slice. The bias is therefore not a perception failure but a
\emph{decision-rule} failure introduced by the anti-bias floor.

\paragraph{The state/trait dichotomy is used to discard valid evidence.}
The chain tags each cue as \texttt{[state]} or \texttt{[trait]}, intended to
separate transient mood from dispositional signal. In practice the verdict
treats the \texttt{[state]} tag as disqualifying. Extraversion record~25 (gold
\textit{high}, predicted \textit{low}) makes the rule explicit:
%
\begin{llmbox}
\textbf{Verdict.}\\
The text \hi{lacks any trait-level evidence} of Extraversion, indicating a low level of this trait. \hi{The presence of state-level cues does not compensate for the absence of stable traits.}
\end{llmbox}
%
The highlighted second sentence states the suppression policy directly: cues
that have been labelled \texttt{state} are given \emph{zero} weight. Since the
upstream tagging routinely files first-person reports of wanting company,
excitement, or social contact under \texttt{state} (they are phrased in the
present tense), the policy systematically zeroes out precisely the language
that signals Extraversion in a diary, driving the $98$-to-$6$ asymmetry on that
trait.

\paragraph{The verdict contradicts its own retrieved-exemplar alignment.}
In the strongest cases the suppression overrides not only the facet scores but
also the retrieval evidence the pipeline was designed to exploit. In
Neuroticism record~3 (gold \textit{high}, predicted \textit{low}) the facet
check scores four facets high and the example-alignment step matches the query
to a \textit{high}-labelled neighbour, yet the verdict still returns
\textit{low}:
%
\begin{llmbox}
\textbf{Facet check.}\\
- anxiety -> \hi{high}; depression -> \hi{high}; self-conscious -> \hi{high}; vulnerability -> \hi{high} (multiple expressions of worry, loneliness, insecurity).\\[4pt]
\textbf{Example alignment.}\\
The test text most closely matches Similar Profile (\hi{label: high}) because it contains multiple expressions of anxiety and emotional distress.\\[4pt]
\textbf{Verdict.}\\
The analysis reveals \hi{no recurring, cross-situational patterns} of emotional instability, leading to a conclusion of low Neuroticism.
\end{llmbox}
%
Here every upstream signal---four high facets and a \textit{high}-labelled
nearest exemplar---points to the correct answer, but the verdict introduces a
\hi{new and stricter criterion} (\enquote{recurring, cross-situational
patterns}) that no single essay can satisfy and overrides them all. This is the
mechanism behind the component ablation in
Table~\ref{tab:pipeline-ablation}: removing the reasoning chain raised macro-F1
on four of five traits precisely because the chain's verdict stage was
\emph{discarding} otherwise-correct retrieval and facet evidence.

\paragraph{The same gate operates on the remaining traits.}
The pattern is not specific to the emotional traits. For Conscientiousness
record~44 the facet check scores \texttt{dutifulness} and \texttt{achievement}
high on the strength of explicit trait-level statements
(\enquote{my purpose here is to stay focused on school\dots}), but the verdict
again invokes the stability gate:
%
\begin{llmbox}
\textbf{Verdict.}\\
The evidence indicates a \hi{lack of stable conscientiousness traits}, with no strong evidence of organization or self-discipline. Therefore, the conclusion is low.
\end{llmbox}
%
and for Agreeableness record~110 the verdict explicitly acknowledges positive
cues before discarding them under the same rule:
%
\begin{llmbox}
\textbf{Verdict.}\\
The analysis indicates a low level of agreeableness due to the \hi{absence of strong trait-level cues}, \hi{despite some situational expressions} of enjoyment and enthusiasm.
\end{llmbox}
%
In both verdicts the highlighted clauses show the model conceding the
high-pole evidence (\enquote{despite some\dots}) and then ruling it out for
failing the trait-level/stability test.

\paragraph{Summary of the mechanism.}
Across all five traits the false negatives share one structure: the evidence
and facet stages correctly surface high-pole cues, and the verdict stage
\emph{suppresses} them by (i) reclassifying present-tense cues as
\enquote{situational}, (ii) giving \texttt{[state]}-tagged evidence zero
weight, and (iii) demanding cross-situational recurrence that a single short
essay cannot demonstrate. The anti-bias floor, intended to stop the model from
defaulting to \textit{high}, therefore over-corrects into a near-blanket
\textit{low} on exactly the traits whose genuine cues are most easily phrased
as transient states---which is why the high-class F1 collapses to $0.166$ on
Neuroticism and $0.358$ on Extraversion in
Table~\ref{tab:model-comparison}, and why deleting the chain entirely
(Table~\ref{tab:pipeline-ablation}) recovers accuracy on most traits.
Calibrating---or removing---this stability gate, rather than changing
retrieval, is thus the primary lever for reducing the bias.

\end{document}
